\section{Introduction}
\label{sec:introduction}

A \(k\)-cap in the projective space \(\PG(d,q)\) over the field of order
\(q\) is a set of \(k\) points, no three collinear.  A line meeting the cap
in two points is a secant, and the cap is complete if every point off it
lies on a secant, that is, if it is maximal under inclusion.  The smallest
size of a complete cap in \(\PG(d,q)\) is denoted \(t_2(d,q)\).  Counting
the points that the \(\binom{k}{2}\) secants can cover gives the classical
lower bound
\begin{equation}
 \binomtwo{k}(q-1)+k\ge\theta_d,
 \qquad
 \theta_j=\frac{q^{j+1}-1}{q-1},
\label{eq:counting-bound}
\end{equation}
so \(t_2(d,q)\ge\sqrt2\,q^{(d-1)/2}\) up to lower-order terms.  In the
plane this bound has been improved by Segre's lemma of
tangents~\cite{Segre1959,Polverino1999}; in \(\PG(3,q)\) a secant-local
count improves the additive constant~\cite{RuddSecantDefects}; for
\(d\ge4\) the literature quotes~\eqref{eq:counting-bound} as the lower
bound, and the exact value of \(t_2(d,q)\) for \(d\ge3\) is known only for
very small \(q\)~\cite{DavydovFainaMarcuginiPambianco2017,
BartoliDavydovKreshchukMarcuginiPambianco2016,BartoliFainaMarcuginiPambianco2017}.
Complete caps are the same objects as parity-check matrices of
quasi-perfect linear codes of minimum distance four and covering radius
two, and \(t_2(d,q)\) is the shortest length of such a code of
codimension \(d+1\)~\cite{DavydovFainaMarcuginiPambianco2017}.

This paper adds one constraint to the counting argument and follows it to
its finite consequences.  The counting bound has a slack, the
\emph{overlap loss}
\begin{equation}
 \Lambda_0=\binomtwo{k}(q-1)-\theta_d+k,
\label{eq:loss}
\end{equation}
which counts, with multiplicity, the points covered by more than one
secant.  The observation is that this slack is not free to distribute:
every hyperplane must carry a prescribed share of it.

\begin{theorem}[Hyperplane loss identity]
\label{thm:hyperplane-loss}
Let \(A\) be a \(k\)-cap in \(\PG(d,q)\), \(d\ge2\), and for a point
\(x\notin A\) let \(r(x)\) be the number of secants of \(A\) through
\(x\).  For every hyperplane \(\Pi\) with \(|A\cap\Pi|=s\),
\begin{equation}
 \sum_{x\in\Pi\setminus A}\bigl(r(x)-1\bigr)=Q(s),
 \qquad
 Q(s)=\binomtwo{k}-\theta_{d-1}+q\binomtwo{s}-(k-2)s.
\label{eq:hyperplane-loss}
\end{equation}
If \(A\) is complete then \(0\le Q(s)\le\Lambda_0\) for every hyperplane
section size \(s\).
\end{theorem}

The lower inequality is the counting bound~\eqref{eq:counting-bound}
applied inside the hyperplane; the upper inequality says that a
hyperplane cannot carry more loss than the whole space has.  Together they
confine the section sizes of a complete cap to the \emph{admissible set}
\begin{equation}
 \cS=\{s:0\le Q(s)\le\Lambda_0\}.
\label{eq:admissible}
\end{equation}
Since \(Q\) is a quadratic in \(s\) with leading coefficient \(q/2\), the
admissible set consists of at most two short runs of integers around the
roots of \(Q\), and when \(\Lambda_0\) is small, which is exactly the
situation at the counting bound, it can have two elements.  A complete
\(31\)-cap in \(\PG(4,7)\), for instance, would have \(\Lambda_0=20\) and
would meet every hyperplane in \(2\) or \(7\) points.

Section sizes of a cap are tied to its parameters by the character
equations, the three moment identities obtained by counting hyperplanes,
point--hyperplane incidences, and pair--hyperplane incidences.  Restricting
those equations to the admissible set gives a finite integer system, and
its infeasibility excludes the cap.  The same equations restricted to the
hyperplanes through a fixed secant \(\ell\) bring in a second invariant,
the number \(T_\ell\) of coplanar four-subsets of \(A\) containing both
endpoints of \(\ell\), which measures how many other secants \(\ell\)
meets.

\begin{theorem}[Secant--hyperplane moments]
\label{thm:secant-hyperplane-moments}
Let \(A\) be a \(k\)-cap in \(\PG(d,q)\), \(d\ge3\), let \(\ell\) be a
secant, and for each hyperplane \(\Pi\supset\ell\) put
\(w_\Pi=|A\cap\Pi|-2\).  Then
\begin{equation}
 \#\{\Pi\supset\ell\}=\theta_{d-2},
 \qquad
 \sum_{\Pi\supset\ell}w_\Pi=(k-2)\theta_{d-3},
 \qquad
 \sum_{\Pi\supset\ell}\binomtwo{w_\Pi}=\binomtwo{k-2}\theta_{d-4}+q^{d-3}\,T_\ell .
\label{eq:moments-intro}
\end{equation}
\end{theorem}

For a complete cap the \(w_\Pi\) lie in \(\cS-2\).  Fixing the count and
the first moment and minimizing the second over that support gives an
unconditional lower bound on \(T_\ell\) for every secant
(Theorem~\ref{thm:local-bounds}), and the secant-local coverage bound
of~\cite{RuddSecantDefects} converts a lower bound on every \(T_\ell\)
into a lower bound on \(\Lambda_0\).  Both routes close finitely.

\begin{theorem}[Exclusions]
\label{thm:exclusions}
There is no complete cap of size \(31\) in \(\PG(4,7)\), \(37\) in
\(\PG(4,8)\), \(97\) in \(\PG(4,16)\), \(293\) in \(\PG(6,8)\), or \(387\)
in \(\PG(6,9)\).  Hence
\[
 t_2(4,7)\ge32,\quad t_2(4,8)\ge38,\quad t_2(4,16)\ge98,\quad
 t_2(6,8)\ge294,\quad t_2(6,9)\ge388 .
\]
\end{theorem}

Each excluded size is the smallest \(k\) satisfying~\eqref{eq:counting-bound},
so each bound improves the counting bound by one.  In four of the five
cases the admissible set has two elements and the global character
equations are already inconsistent over it; the fifth, \(\PG(6,8)\), has
five admissible sizes separated by a gap of eleven around the mean, and is
excluded by the coverage route.  A deterministic exact-arithmetic scan
over \(d\le6\) and small \(q\) (Section~\ref{sec:exclusions}) finds no
other gain over the counting bound in dimension at least four, and never a
gain of more than one.  Section~\ref{sec:limits} explains why: one step
above the counting bound \(\Lambda_0\) grows by about \(kq\), the
admissible set becomes an interval, and every test is satisfied.  The value
of Theorem~\ref{thm:exclusions} is therefore the mechanism rather than the
numbers, which remain far from the smallest known complete caps
(\(56\) points in \(\PG(4,7)\)~\cite{BartoliDavydovKreshchukMarcuginiPambianco2016}).

The second contribution is negative.  The companion
paper~\cite{RuddSecantDefects} asks whether a positive proportion of the
secants of a complete cap in \(\PG(d,q)\), \(d\ge4\), must carry
\(T_\ell\ge cq\) collisions, which would improve the counting bound at the
scale \(q^{(d-3)/2}\).  Proposition~\ref{prop:degree-countermodel}
constructs formal collision data satisfying every constraint that the
coplanar-quadruple counts, the secant--point incidences, and the
plane-pencil identities impose, in which only \(o\bigl(\binom{k}{2}\bigr)\)
secants meet any other.  Any proof of such a proportion must use an
ambient constraint; Theorem~\ref{thm:hyperplane-loss} is one.

\subsection*{Related work}

The global counting bound~\eqref{eq:counting-bound} and the dictionary
between complete caps and quasi-perfect codes are standard; the papers of
Bartoli, Davydov, Faina, Kreshchuk, Marcugini, and Pambianco cited above
collect the known sizes and constructions of small complete caps and
quote \(\sqrt2\,q^{(d-1)/2}\) as the lower bound for \(d\ge3\).  The
character equations for hyperplane sections are classical; in coding
terms they are the first two power moment identities of
Pless~\cite{Pless1963} for the projective code generated by the cap, and
Theorem~\ref{thm:secant-hyperplane-moments} is the same identity for the
code shortened at two coordinates.  Counting cap points over the
hyperplanes through a fixed subspace is a standard device, used for
instance by Thas over the hyperplanes through a plane to bound the
largest caps of \(\PG(n,q)\), \(q\) even~\cite{Thas2017}.  Divisibility
conditions on hyperplane intersection numbers are the subject of the
theory of divisible codes~\cite{Kurz2021}; the divisibility used here is
imposed by completeness rather than assumed.  Wehlau classifies the
complete caps of \(\PG(n,2)\) that miss a codimension-two subspace
through their hyperplane intersection sizes~\cite{Wehlau2004}, the
nearest instance we know of constraining a complete cap by its hyperplane
sections; the mechanism there is the binary structure, not a loss count.
Pairs with \(T_\ell=0\) are the free pairs of Farr and
Lison\v{e}k~\cite{FarrLisonek2006a,FarrLisonek2006}, and caps in which
every pair is free are the \(4\)-general sets of
Pavese~\cite{Pavese2025}; both notions reappear in
Section~\ref{sec:limits}.  The overlap loss \(\Lambda_0\) and the
secant-local statistics \(T_\ell\) and \(\varphi_m\) are taken
from~\cite{RuddSecantDefects}, and the integer-envelope viewpoint from
\cite{RuddIntegralSecant}; the present paper differs from both in that its
degree sequences are hyperplane section sizes rather than secant
occupancies or point--block incidences.  The lower half of
Theorem~\ref{thm:hyperplane-loss} is a one-line consequence of
completeness that we expect to be folklore, though we did not find it
stated; we found no earlier statement of the upper half, of the
admissible set, or of any exclusion of a counting-bound size for a
complete cap in \(\PG(d,q)\) with \(d\ge4\).  The searches and sources
behind that sentence are recorded with the paper's repository.

\subsection*{Organization}

Section~\ref{sec:preliminaries} fixes notation and recalls the pencil and
character identities.  Sections~\ref{sec:hyperplane-loss}
and~\ref{sec:moments} prove Theorems~\ref{thm:hyperplane-loss}
and~\ref{thm:secant-hyperplane-moments}, including the version for caps
complete outside a prescribed set of holes.
Section~\ref{sec:feasibility} derives the integer feasibility conditions
and the local bounds on \(T_\ell\), Section~\ref{sec:coverage} converts
them into coverage bounds, and Section~\ref{sec:exclusions} proves
Theorem~\ref{thm:exclusions} and describes the scan.
Section~\ref{sec:limits} gives the formal countermodel and the reasons the
method stops one step above the counting bound.
